% The classical model as a projection --- plan \S4.4.
% Sources: new_notes3 Part D \S20; new_notes Part II \S10--\S12; new_notes2
% \S8, plus the overlays migrated from draft \S3.4.  The p_eff example
% values and the trade-off inputs were re-verified numerically on
% 2026-08-07.  The generation kernel of \S4.2 is the licensed addition of
% plan \S9.

\section{The classical model as a projection}
\label{p:sec:projection}

The renewal system does not merely disagree with classical BMVR in
transients; it contains classical BMVR as an exact special case, and the
relationship is precise enough to be turned into predictions.

\subsection{Exact reduction in the exponential phase}

Suppose the incidence grows exponentially, $i(t)=i_0\ee^{rt}$, at times late
enough that the lower limits of the convolutions are irrelevant. Then
\begin{equation}
\Icell(t)=i_0\ee^{rt}\Lap{\Ifix}(r),
\qquad
\Vfree'(t)=i_0\ee^{rt}\,\delta\Lap{K}(r)-c\,\Vfree(t),
\label{p:eq:expphase}
\end{equation}
where $\Lap{f}(r)=\int_0^\infty\ee^{-r\alpha}f(\alpha)\,\dd\alpha$. A
classical model driven by the same incidence has
$\Icell(t)=i_0\ee^{rt}/(r+d_{\Icell})$ and the eigenvalue relation
$(r+d_{\Icell})(r+c)=\gamma Tp$; matching the two exponential solutions
fixes
\begin{equation}
p_{\rm eff}(r)=\frac{\delta\,\Lap{K}(r)}{\Lap{\Ifix}(r)},
\qquad
d_{\Icell,{\rm eff}}(r)=\frac{1}{\Lap{\Ifix}(r)}-r.
\label{p:eq:peff}
\end{equation}

The reduction is exact, and it is the central claim of the chapter.

\begin{theorem}[Exponential-phase equivalence]
\label{p:thm:projection}
Fix a growth rate $r>-\theta$. The constant pair
$\bigl(p_{\rm eff}(r),d_{\Icell,{\rm eff}}(r)\bigr)$ of \cref{p:eq:peff} is
the unique one for which classical BMVR reproduces the renewal system's
$\ee^{rt}$ solution exactly, in both $\Icell$ and $\Vfree$; and substituting
it into the classical eigenvalue relation returns
\begin{equation}
r+c=\gamma T\,\delta\Lap{K}(r),
\label{p:eq:char}
\end{equation}
precisely the characteristic equation of the renewal system. What no
constant pair reproduces is a transient, which samples a continuum of growth
rates at once.
\end{theorem}

\begin{proof}
Uniqueness and existence are \cref{p:eq:expphase} read as two equations in
two unknowns: matching $\Icell$ forces
$1/(r+d_{\Icell})=\Lap{\Ifix}(r)$, which is $d_{\Icell,{\rm eff}}$, and
matching the release term then forces
$p\,\Lap{\Ifix}(r)=\delta\Lap{K}(r)$, which is $p_{\rm eff}$. Both
transforms are finite and strictly positive for $r>-\theta$, so the pair is
well defined. Substituting into $(r+d_{\Icell})(r+c)=\gamma Tp$ gives
$(r+c)/\Lap{\Ifix}(r)=\gamma T\,\delta\Lap{K}(r)/\Lap{\Ifix}(r)$, that is
\cref{p:eq:char}. That no single pair serves a transient follows from
\cref{p:prop:peff-monotone} below: $p_{\rm eff}$ is not constant in $r$.
\end{proof}

\noindent
So the classical model is not a rival to the renewal system but a projection
of it, exact on one exponential phase at a time. Everything else in this
section is a consequence of reading \cref{p:eq:peff} as a map.

The two ends of the map are explicit. Using $\int_0^\infty\delta K=V_\infty$ (from $V'=\delta K$) and
$\int_0^\infty\Ifix=\mean{T_{\rm prod}}$ (\cref{p:eq:lifetime}),
\begin{equation}
r=0:\qquad
p_{\rm eff}(0)=\frac{V_\infty}{\mean{T_{\rm prod}}}
=\frac{\lambda\,V_\infty}{\log\bigl(a/(a-1)\bigr)},
\qquad
d_{\Icell,{\rm eff}}(0)=\frac{\lambda}{\log\bigl(a/(a-1)\bigr)};
\label{p:eq:peff0}
\end{equation}
the $r=0$ value coincides with naive moment-matching, lifetime output
divided by lifetime --- a consistency check, and the matching used in the
overlays below. At the other end,
\begin{equation}
r\to\infty:\qquad
p_{\rm eff}(r)\to\delta,
\qquad
d_{\Icell,{\rm eff}}(r)\to\mu+\delta,
\label{p:eq:peffinf}
\end{equation}
the \emph{young-cell} limit: in a fast-growing epidemic almost every
infected cell was infected moments ago, holds one particle, bursts at rate
$\delta$, and releases one.

Calling these the two ends understates the map, because $r=0$ is not an end.
Both transforms converge for every $r>-\theta$, and negative $r$ is not a
fiction: it is a decaying infection, the case of
\cref{p:fig:overlay-V}(b,d). The genuine left-hand endpoint is the mirror of
the young-cell limit.

\begin{proposition}[The old-cell limit]
\label{p:prop:oldcell}
The release flux per surviving cell tends to a finite limit,
\begin{equation}
\lim_{\alpha\to\infty}\frac{g(\alpha)}{S(\alpha)}
=\frac{a\lrbs{2\lambda(1-b)+\delta}}{a-1}
=\frac{\delta\,a(a+1)}{(a-1)^2}
=\delta\,\mean{X^2}_{\rm QS},
\label{p:eq:oldcell}
\end{equation}
the burst rate times the second moment of the quasi-stationary load, and
consequently $p_{\rm eff}(r)\to\delta\,\mean{X^2}_{\rm QS}$ as
$r\downarrow-\theta$.
\end{proposition}

\begin{proof}
Write $v=\ee^{-\theta\alpha}$. As $\alpha\to\infty$, $v\to0$ and
\cref{p:eq:IDIhat,p:eq:J} give $\Ifix\sim(a-b)^2v/(Aa)$ and $I\to b$, so by
\cref{p:eq:K} $g=2\lambda(\kappa-I)J\sim2\lambda(\kappa-b)(a-b)^2v/A^2$. The
ratio of the two is $2\lambda(\kappa-b)a/A$, which is the first equality
after $2\lambda(\kappa-b)=2\lambda(1-b)+\delta$. For the second, the
identity $AB=\delta/\lambda$ of \cref{p:eq:AB} reads
$\lambda(1-b)=\delta/(a-1)$, whence
$a[2\delta/(a-1)+\delta]/(a-1)=\delta a(a+1)/(a-1)^2$. That last expression
is $\delta\mean{X^2}_{\rm QS}$ because the quasi-stationary load is
geometric on $\{1,2,\ldots\}$ with ratio $1/a$ (\cref{p:eq:burstlaw}), whose
second moment is $a(a+1)/(a-1)^2$. Finally, both $\Lap{\Ifix}$ and
$\delta\Lap K$ diverge like $1/(r+\theta)$ as $r\downarrow-\theta$, with
coefficients in the ratio just computed, so their quotient tends to it.
\end{proof}

Two things are worth saying about \cref{p:eq:oldcell}. It repairs the failed
proposal of \cref{p:sec:obstruction}, which offered the quasi-stationary
\emph{mean} $\mean{X}_{\rm QS}$ as a release rate and was wrong twice over:
a count was used where a rate was needed, and the first moment where the
second is called for. Both corrections appear in \cref{p:eq:oldcell}
together --- $\delta$ supplies the missing inverse time, and the
load-proportional burst hazard weights by $X$ a second time. The naive
proposal was, in that sense, an old-cell limit with the wrong moment and no
clock. And it gives the absorbing process the mature-cell limit that
\cref{p:sec:reset} gives the resetting one, in the same form
$p_\infty=\delta\,\mathbb E_\pi[X^2]$ with $\pi$ the quasi-stationary law.

Between the endpoints the map is monotone, and the reason is a single
property of the kernels.

\begin{proposition}[Monotonicity of the effective release rate]
\label{p:prop:peff-monotone}
Write $\rho(\alpha)=g(\alpha)/S(\alpha)$ for the expected release flux per
still-productive cell, so that $\rho(0)=\delta$ and
$\rho(\infty)=\delta\mean{X^2}_{\rm QS}$. If $\rho$ is non-decreasing then
$p_{\rm eff}$ is non-increasing on $(-\theta,\infty)$, strictly so if $\rho$
is strictly increasing.
\end{proposition}

\begin{proof}
For $r>-\theta$ let $\pi_{f,r}$ be the probability density
$\alpha\mapsto\ee^{-r\alpha}f(\alpha)/\Lap f(r)$. Differentiating under the
integral,
$\dd\log\Lap f(r)/\dd r=-\mean{\alpha}_{\pi_{f,r}}$, so
\begin{equation}
\frac{\dd}{\dd r}\log p_{\rm eff}(r)
=\mean{\alpha}_{\pi_{S,r}}-\mean{\alpha}_{\pi_{g,r}}.
\label{p:eq:peff-deriv}
\end{equation}
The two densities have ratio
$\pi_{g,r}(\alpha)/\pi_{S,r}(\alpha)\propto\rho(\alpha)$, independently of
$r$. If $\rho$ is non-decreasing, $\pi_{g,r}$ dominates $\pi_{S,r}$ in the
monotone likelihood ratio order, hence in the usual stochastic order, hence
in mean; so \cref{p:eq:peff-deriv} is non-positive.
\end{proof}

\noindent
The hypothesis is a statement about one closed-form function of one
variable: substituting $v=\ee^{-\theta\alpha}$ into
\cref{p:eq:IDIhat,p:eq:J,p:eq:K} gives
\begin{equation}
\rho=\frac{2\lambda\lrbs{B(\kappa-a)v+A(\kappa-b)}(a-bv)}{(Bv+A)^2},
\label{p:eq:rho-closed}
\end{equation}
which must be shown decreasing in $v\in(0,1]$. That reduces the question to
a finite algebraic one; it is not settled here, but $\rho$ was found
increasing at every one of $4000$ triples sampled over
$(\mu,\delta)\in[10^{-3},20]^2$ at $\lambda=1$, and $p_{\rm eff}$ decreasing
at every one (test~M of \cref{p:app:verification}).

The size of the effect is what matters. At
$(\lambda,\mu,\delta)=(1,0,0.1)$, $p_{\rm eff}(0)\approx4.59$ against
$p_{\rm eff}(\infty)=0.1$ --- the same cell, two orders of magnitude apart,
with nothing separating the two figures but how fast the epidemic around it
is growing. Over the whole map the span is larger, and exact: by
\cref{p:eq:peffinf,p:eq:oldcell} the ratio of the two endpoints is
\begin{equation}
\frac{\delta\,\mean{X^2}_{\rm QS}}{\delta}=\mean{X^2}_{\rm QS}
=\frac{a(a+1)}{(a-1)^2},
\label{p:eq:peff-span}
\end{equation}
a pure number fixed by the intracellular rates alone. It is $231$ at
$(1,0,0.1)$ and $30.5$ at $(1,0.9,0.1)$: the full dynamic range of a fitted
release rate is the second moment of the quasi-stationary load.

\begin{figure}[tbp]
\centering
\includegraphics[width=0.94\textwidth]{peff_dr_curves.pdf}
\caption{The effective release rate (a) and removal rate (b) of
\cref{p:eq:peff} against the epidemic growth rate $r$, for three parameter
sets, over the whole domain $r>-\theta$ on which the map is defined. Open
circles mark the $r=0$ values \cref{p:eq:peff0}; dotted segments mark the
limits each curve runs toward. Note that $r=0$ is an interior point, not an
end: the left-hand endpoint is the old-cell limit $\delta\,\mathbb{E}_{\rm
QS}[X^2]$ and the right-hand one the young-cell limit $\delta$. Read the
vertical extent of each curve as the size of the error made by treating $p$
as a constant of the pathogen rather than of the epidemic's phase.}
\label{p:fig:peff}
\end{figure}

The closed hypergeometric forms of $\Lap{\Ifix}$ and $\delta\Lap{K}$ are in
\cref{p:app:transforms}; the curves above are computed by direct quadrature.

\subsection{The generation kernel}
\label{p:sec:generation-kernel}

One object lies behind both \cref{p:eq:peff} and the comparison of
\cref{p:sec:flooding}. Define the cell-to-cell \emph{generation kernel}
\begin{equation}
\mathcal A(\alpha)
=\gamma T\,\bigl(g*\ee^{-c\,\cdot}\bigr)(\alpha)
=\gamma T\!\int_0^\alpha g(u)\,\ee^{-c(\alpha-u)}\,\dd u,
\label{p:eq:genkernel}
\end{equation}
the rate at which a cell of age $\alpha$ generates new infected cells: it
releases at rate $g(u)$ at age $u$, a particle released then is still free a
time $\alpha-u$ later with probability $\ee^{-c(\alpha-u)}$, and infects at
rate $\gamma T$. Its transform is
\begin{equation}
\Lap{\mathcal A}(r)=\frac{\gamma T\,\Lap{g}(r)}{r+c},
\label{p:eq:Atilde}
\end{equation}
so the characteristic equation \cref{p:eq:char} is exactly
$\Lap{\mathcal A}(r)=1$: the Euler--Lotka equation of a
one-type age-structured population, with the infected cell as the
individual and $\mathcal A$ as its mean reproduction measure
\cite{diekmann1990definition,metz1986dynamics}.
% NEEDS-REF: Crump--Mode--Jagers general branching processes, the umbrella
% object of which A is the mean reproduction measure and r the Malthusian
% parameter

Three things follow, each of which the chapter already contains in another
form. First, the total mass of $\mathcal A$ is the basic reproduction
number,
\begin{equation}
R_0=\int_0^\infty\mathcal A(\alpha)\,\dd\alpha
=\gamma T\int_0^\infty g\;\cdot\;\int_0^\infty\ee^{-cu}\,\dd u
=\frac{\gamma T\,V_\infty}{c},
\label{p:eq:R0gen}
\end{equation}
which depends on $g$ only through its integral. Second, the normalised
kernel $\mathcal A(\alpha)/R_0$ is the \emph{generation-time density}, the
law of the age at which a cell produces a randomly chosen one of its
offspring cells; so the mechanism named in \cref{p:sec:tradeoff} is a
definition rather than an interpretation, and the numbers that go with it
follow from $g$ alone. By
\cref{p:eq:genkernel} the mean generation time is the mean release age under
$g/V_\infty$ plus the mean extracellular wait $1/c$: for the budding
comparator matched at $r=0$ the release age is exponential with mean
$\mean{T_{\rm prod}}=2.398$, giving $2.398+1=3.398$ at $c=1$; for the
bursting cell it is $3.089$, giving $3.089+1=4.089$. Third, since $r$ solves
$\Lap{\mathcal A}(r)=1$ while $R_0=\Lap{\mathcal A}(0)$, threshold and growth
rate are different functionals of one kernel --- the standard
$R_0$-versus-$r$ distinction \cite{wallinga2007generation}, here with both
sides in closed form.

\subsection{What the projection loses}
\label{p:sec:overlays}

To see what age structure does to an epidemic, integrate both systems side
by side. In \cref{p:fig:overlay-V,p:fig:growth-phase} the solid curve is the
renewal system and the dashed one is classical BMVR with the fairest
constant match, at the $r=0$ end of \cref{p:eq:peff},
\begin{equation}
p=p_{\rm eff}(0)=\frac{V_\infty}{\mean{T_{\rm prod}}}
=\frac{\lambda V_\infty}{\log\bigl(a/(a-1)\bigr)},
\qquad
d_{\Icell}=d_{\Icell,{\rm eff}}(0)=\frac{\lambda}{\log\bigl(a/(a-1)\bigr)},
\label{p:eq:r0match}
\end{equation}
so that total lifetime output and mean productive lifetime agree. The
comparator is the budding cell of \cref{p:sec:classical}, its constants now
chosen by the map rather than by hand.

\begin{figure}[tbp]
\centering
\includegraphics[width=0.94\textwidth]{overlay_V.pdf}
\caption{Both state variables, both sides of the threshold. Free particles
$\Vfree(t)$ (top row) and infected cells $\Icell(t)$ (bottom) under the
renewal system (solid) and under classical BMVR matched at $r=0$ (dashed),
for a supercritical infection (left, $R_0=1.76$) and a subcritical one
(right, $R_0=0.28$). The matched model gets the order of magnitude and not
the shape; and the discrepancy is not confined to the release equation ---
age structure moves the cell trajectory too, which is why matching $p$ alone
cannot repair it. Parameters: $(\lambda,\mu,\delta)=(1,0,0.1)$ throughout,
with $\gamma T=0.04$, $c=0.25$ on the left and $\gamma T=0.01$, $c=0.4$ on
the right.}
\label{p:fig:overlay-V}
\end{figure}

The reading is consistent across panels. The matched model misses three
things: the early lag, because young cells dominate and release is nearer
the young-cell limit $p\to\delta$ than $p_{\rm eff}(0)$; the timing of the
peak; and the late growth rate, since generation-time structure moves the
operative release rate away from its $r=0$ value as the epidemic grows. The relative
free-particle discrepancy,
\begin{equation}
\frac{\bigl|\Vfree^{\rm ren}(t)-\Vfree^{\rm cl}(t)\bigr|}
{\max_s \Vfree^{\rm ren}(s)},
\label{p:eq:reldiff}
\end{equation}
is often order one or larger by the end of the plotted window in
supercritical regimes, and worst when intracellular death is high, where
$V_\infty$ is small while the age structure of release stays strongly
non-exponential. Discrepancies grow with $R_0$ and with that age structure
--- slow bursting, non-zero $\mu$.

\begin{figure}[tbp]
\centering
\includegraphics[width=0.94\textwidth]{overlay_growth_phase.pdf}
\caption{Why no single constant pair will do. One renewal trajectory
(solid) against two classical parameter choices: the $r=0$ match (dashed)
and the young-cell limit $p=\delta$, $d_{\Icell}=\mu+\delta$ (dash-dot).
(a) The first six time units, where the $r=0$ match overshoots by more than
an order of magnitude while the young-cell parameters are the closer of the
two. (b) The full window, where the ordering has reversed and the young-cell
model has collapsed. Each choice is correct at one growth rate and the
trajectory passes through many, exactly as \cref{p:eq:peff} predicts. Both
panels logarithmic in $\Vfree$. Parameters:
$(\lambda,\mu,\delta)=(1,0,0.1)$, $\gamma T=0.025$, $c=0.2$, $R_0=1.38$.}
\label{p:fig:growth-phase}
\end{figure}

These are the facts the reduction predicts: a constant pair reproduces one
exponential growth rate exactly, and a transient sampling a continuum of
ages is beyond it.

\subsection{The invariance of \texorpdfstring{$R_0$}{R0}}

Setting $r=0$ in \cref{p:eq:char} gives $\gamma T\,\delta\Lap{K}(0)=c$,
that is $\gamma TV_\infty=c$, so the renewal system's basic reproduction
number is
\begin{equation}
R_0^{\mathrm{ODE}}=\frac{\gamma T\,V_\infty}{c}.
\label{p:eq:R0ODE}
\end{equation}
This agrees with \cref{p:eq:R0gen} \cite{diekmann1990definition}.

\begin{theorem}[$R_0$ is unchanged by bursting]
\label{p:thm:R0inv}
At fixed total per-cell output $V_\infty$, replacing constant budding by
bursting leaves $R_0$ exactly unchanged. Bursting alters the
generation-time distribution --- hence the growth rate $r$, hence the
extinction probability, hence the stochastic early phase --- but not the
threshold itself.
\end{theorem}

\begin{proof}
By \cref{p:eq:R0gen}, $R_0=\int_0^\infty\mathcal A$ depends on the release
kernel only through $\int_0^\infty g=V_\infty$. Replacing bursting by
budding at matched yield replaces $g$ by another kernel with the same
integral, and leaves $R_0$ where it was. The generation-time density
$\mathcal A/R_0$, by contrast, depends on the whole of $g$, and it is
through that density that $r$ is determined by $\Lap{\mathcal A}(r)=1$.
\end{proof}

Two remarks attach to the theorem. The first is evidential: the numerical
checks of \cref{p:fig:EF} put the late growth rate at the characteristic
root $r_*$ of \cref{p:eq:char} to about $10^{-7}$ across five scenarios, and
trajectories decay for $R_0<1$ and grow for $R_0>1$, with the transition
resolved already at $R_0=1.05$. The second is a caution about assumption
(A2): \cref{p:eq:R0ODE} is the \emph{continuum} reproduction number, in which one
virion contributes mean $\gamma T/c$ new infections, whereas a discrete
branching process absorbing virions on infection carries the mean
$m=qV_\infty$ with $q=\gamma T/(\gamma T+c)$. The two agree when $\gamma
T\ll c$, and it is $m$ that governs \cref{p:sec:flooding}.

\subsection{Identifiability: a negative result}
\label{p:sec:identifiability}

How much of the intracellular process can a population-level fit recover? At
any fixed growth rate $r$, the map
\begin{equation}
(\lambda,\mu,\delta)\ \longmapsto\ \lrb{p_{\rm eff}(r),\,d_{\Icell,{\rm eff}}(r)}
\label{p:eq:identmap}
\end{equation}
sends three rates to two numbers, and is therefore not invertible
(\cref{p:fig:identifiability}). This is a structural identifiability
statement of the kind Miao et al.\ \cite{miao2011identifiability} treat for
viral dynamics generally, and it is negative for the same reason: the
observable is a projection, and the projection has a kernel. It also closes
the ambition of \cref{p:sec:obstruction}, which hoped that a good estimate
of $p$ would give access to the rates making up the quasi-stationary mean:
it does not, and neither does $(p,d_{\Icell})$ together.

\begin{figure}[tbp]
\centering
\includegraphics[width=0.94\textwidth]{N4b_2_identifiability_levels.pdf}
\caption{Population observables leave level sets rather than identifying the
intracellular rates. With $c=1$ and $\gamma T=1$, the interior flooding
point $(\lambda,\mu,\delta)=(1,0.5,1/3)$ has $R_0=2$ and $r_\star=0.294$.
(a) One population observable leaves a level curve: in the slice $\lambda=1$ the contours $R_0=2$ and $r=r_\star$ are
distinct; alternatives A and C share $R_0$ with the reference but not its
growth rate. (b) With $\lambda$ free, a continuous family of triples shares
both. Joint knowledge of threshold and exponential growth therefore does not
identify $(\lambda,\mu,\delta)$.}
\label{p:fig:identifiability}
\end{figure}

A third observable is required; in order of experimental accessibility:
\begin{enumerate}
\item the burst-size dispersion --- the geometric ratio $1/a$ fixes one
combination of the rates directly (the burst-size law of \ChDist{});
\item the fraction of cells that clear without bursting, which is $b$;
\item the shape --- not merely the mean --- of the release-time density
$\delta J$ (\ChDist{}).
\end{enumerate}
Any one of these, with $(p,d_{\Icell})$, over-determines
$(\lambda,\mu,\delta)$ and converts the fit into a consistency test of the
intracellular model. The limitation is not one of fitting procedure but of
the information content of classical BMVR itself.

\subsection{A prediction for fitting practice}

The fitting consequences --- acute-phase versus set-point estimates of $p$,
and the eclipse-like lag visible in
\cref{p:fig:overlay-V,p:fig:growth-phase} --- are taken up in
\cref{p:sec:fitting-consequences}. The lag \emph{is} the eclipse: it
emerges from the intracellular dynamics with no extra parameter.

\FloatBarrier
